% Transcription — fonds Alexandre Grothendieck, shelfmark 159, batch 2 (pages 21-33).
% Established from the facsimile archives/batches/159/batch-02.pdf (a mirror of
% https://grothendieck.umontpellier.fr/159.pdf), per the skill
% transcribe-grothendieck. The folder is thirty-three pages; this is its second
% and last batch, the thirteen leaves 21-33. The facsimile has no cover sheet:
% its page 1 is the archivists' page 21, checked against the pencilled numbers
% bottom-left.
%
% Page 23 is a blank leaf and is absent; every other leaf of the batch is
% mathematics and is transcribed. There is no cover and no pagination of his
% own. The run is not continuous: page 22 is three lines finishing page 21,
% and pages 24, 25-26, 27-28/30, 29, 31, 32-33 are so many restarts of the
% same material, in the manner of the folder.
%
% Content. Pages 21-22 close the argument begun earlier in the folder: for E
% of rank 2 with det E = ω_X, det E ≃ det F ⊗ L^⊗2 forces deg L = g-1, and the
% composite ω ↪ E ≃ F⊗L ↠ L is a section of L⊗ω^{-1} of negative degree, hence
% zero — whence a contradiction. The theorem drawn from it: for X projective
% smooth of genus ≥ 2, no projective connection on X lets the structure group
% of the associated principal bundle (group H = PGL(1)_S in his notation of
% page 1) reduce to a proper algebraic subgroup — neither a Borel nor the
% normaliser of a maximal torus, since the double cover would reduce it to the
% torus. Over C this says the monodromy π_1(X^an,x) → H(C) has Zariski-dense
% image.
%
% Pages 24-33 are a second subject: the formal theory. P is an adically
% complete augmented O_X-Algebra with J/J² = Ω invertible and Ω ≃ Ω¹_{X/S}, C
% a connection on P relative to S compatible with the algebra structure, and
% Ŝ = Sym^(Ω) the model. The theorem (page 25) asserts a unique isomorphism
% u : Ŝ → P of augmented algebras with u(C) = ĉ; a) the datum of u to order 2
% already determines (P,σ) — i.e. the extension 0 → Ω → E → O_S → 0 — up to
% unique isomorphism, b) any u_3 to order 3 extending it extends uniquely to
% u. Page 26 turns this into a bijection between the solutions, the pseudo-u_3
% prolonging u_2 (a torsor under Γ(X, Ω^⊗2)) and Conn*(P). Pages 27, 28 and 30
% are the explicit computation of C(ω_0^u) in Ŝ, where the raised u denotes the
% factor u = 1 + Σ_{i≥2} α_i ω_0^i; the outcome is that the change of parameter
% contributes (i+1)α_i ω_0^i ϖ_0 and that regularity forces (n+1)α_n = 0.
% Page 29 sets G (automorphisms of Ŝ trivial mod J³, so G/G_0 ≃ Γ(Ω^⊗2)),
% G(c) = {g : gĉ of the form ĉ'}, and proves G_0 × G(c) → G bijective. Page 31
% strips that proof to an abstract lemma on a group G ⊃ G_0 acting on C ⊃ C_0.
% Pages 32-33 restate the setting once more and define a connection to be
% "régulière" when ϖ_0 : Ω → Ω¹_{X/S} is an isomorphism, closing on
% G(c) ≃ G/G_0 ≃ Γ(X, Ω^⊗2).
%
% Reading hazards fixed once here. The coefficients of C(ω_0) are written with
% a single looped stroke and are set ϖ_i throughout. The raised u of ω_0^u is
% a factor, not an exponent — page 30 writes the same thing out in full. J, Ĵ,
% Ŝ, ĉ and the hats are his. Underlining is set in bold or by \underline as he
% uses it.
%
% Readings flagged and not settled: on page 21, « lorsque », « normalisateur »,
% « passage au revêtement double », the group name read PGL(1)_S, and a single
% symbol repeated twice before « fibrés » which is left \ill{}; on page 25, the
% parenthesis after the J-adic filtration, the sign between Ω and (J/J³), and
% the phrase introducing the category of extensions; on page 26, item c) in
% its entirety; on page 28, the index letter of α_n, drawn like a μ; on page
% 32, the letter after « sép. complète pour »; on page 33, « régulière », which
% is however settled by the gloss ϖ_0 : Ω ≃ Ω¹_{X/S} two lines below. The
% lower halves of pages 26, 27 and 30 are scratch calculation, largely struck,
% and carry many more marks than the statements do.
%
% Pass: Opus 5 (claude-opus-5), 2026-09-19 — first pass, unchecked against the pages by a human.
\documentclass[11pt,a4paper]{article}
% Le préambule commun à toutes les transcriptions du fonds Grothendieck.
%
% Il tient en une idée : **l'incertitude se compose, elle ne se résout pas.**
% Une transcription qui lisse un mot douteux a détruit la seule chose pour
% laquelle on la faisait — un lecteur doit pouvoir savoir, à chaque endroit,
% ce qui était sur le feuillet et ce qui a été ajouté. Les six macros qui
% suivent sont donc l'appareil critique, et non de la décoration.
%
% Les mêmes commandes sont reconnues par `scripts/render.mjs`, qui produit la
% vue de lecture du site. Une macro ajoutée ici doit l'être aussi là-bas,
% faute de quoi le rendu échouera bruyamment — ce qui est voulu : un rendu
% silencieusement faux est pire qu'un rendu absent.

\ProvidesPackage{grothendieck}[2026/08/08 Transcriptions du fonds Grothendieck]

\usepackage{amsmath,amssymb,amsthm}
\usepackage{mathrsfs}
\usepackage{stmaryrd}
% Les diagrammes commutatifs sont partout dans ce fonds : c'est la langue
% même de la géométrie algébrique de Grothendieck, et une transcription qui
% les rendrait en prose ne transcrirait rien.
\usepackage{tikz-cd}
% Français par défaut — c'est la langue des feuillets — mais l'anglais est
% chargé aussi : la traduction et le résumé sont des documents anglais, et la
% typographie française (espaces avant les deux-points) y serait une faute.
% Ces documents basculent par \selectlanguage{english} en tête de corps.
\usepackage[english,french]{babel}
\usepackage{fontspec}
\usepackage{geometry}
\geometry{a4paper,margin=25mm,marginparwidth=28mm}
\usepackage{marginnote}
\usepackage{xcolor}
% ulem, not soul. `soul`'s \st and \ul are built on a character-by-character
% scan that XeLaTeX's font machinery breaks: the document fails with an
% undefined control sequence rather than degrading. ulem does the same two jobs
% and is Unicode-engine safe. `normalem` stops it hijacking \emph, which the
% transcriptions use for Grothendieck's own emphasis.
\usepackage[normalem]{ulem}
\usepackage{hyperref}
\hypersetup{colorlinks=true,linkcolor=black,urlcolor=black,pdfborder={0 0 0}}

% --- Métadonnées du lot -----------------------------------------------------
% Reprises telles quelles par le script de rendu, qui les affiche en tête de
% la vue de lecture. Elles fixent ce que le fichier prétend couvrir : sans
% elles, une transcription flotte sans point d'attache dans un fonds de seize
% mille feuillets.

\newcommand{\folder}[1]{\gdef\grothFolder{#1}}
\newcommand{\batch}[1]{\gdef\grothBatch{#1}}
\newcommand{\leaves}[2]{\gdef\grothFirst{#1}\gdef\grothLast{#2}}
\newcommand{\foldertitle}[1]{\gdef\grothFoldertitle{#1}}
\gdef\grothFolder{?}\gdef\grothBatch{?}\gdef\grothFirst{?}\gdef\grothLast{?}\gdef\grothFoldertitle{}

% --- Le résumé --------------------------------------------------------------
%
% Il ouvre la lecture modernisée, sous le titre « Résumé », et il est composé
% en petit. Ce n'est pas un ornement : c'est le seul passage du document qui
% ne soit pas des mathématiques, et un lecteur doit pouvoir le reconnaître
% avant de l'avoir lu — puis le sauter s'il connaît déjà le sujet, ou s'y
% arrêter s'il ne le connaît pas. Le filet qui le suit marque la reprise du
% corps.

\newenvironment{resume}
  {\small\setlength{\parskip}{0.45em}}
  {\par\medskip\hrule\medskip}

% --- Mots-clés ---------------------------------------------------------------
%
% En anglais, en clôture du résumé de la lecture modernisée : le vocabulaire
% moderne sous lequel chercher ce que les feuillets construisent. Le manifeste
% du site (`npm run manifest`) les extrait pour étiqueter la cote entière —
% cette ligne est la source unique des tags, il n'y a pas de fichier à côté.
\newcommand{\keywords}[1]{\par\smallskip\noindent
  {\footnotesize\textsc{Keywords} --- \textit{#1}}\par}

% --- L'appareil critique ----------------------------------------------------

% Le numéro de feuillet, en marge. C'est l'ancre qui permet de régler un
% désaccord sur une formule en regardant *un* feuillet plutôt qu'une chemise
% de six cents. Le site s'en sert aussi pour tourner les pages du fac-similé
% quand on fait défiler la transcription.
\newcommand{\leaf}[1]{%
  \marginnote{\footnotesize\color{gray!70}\textsf{#1}}%
  \label{leaf:#1}\ignorespaces}

% Illisible : on ne devine pas. Un mot inventé qui se lit comme les autres est
% le pire résultat possible.
\newcommand{\ill}{\textcolor{red!60!black}{[\,\ldots\,]}}

% Lecture douteuse : le mot est donné, et signalé comme incertain.
\newcommand{\uncertain}[1]{\uline{#1}}

% Ajout de l'éditeur — une lettre manquante, un mot sous-entendu.
\newcommand{\add}[1]{\textcolor{blue!50!black}{[#1]}}

% Biffé par Grothendieck lui-même. Ce qu'il a rayé fait partie du document :
% une rature dit souvent où la pensée a changé de direction.
\newcommand{\struck}[1]{\sout{#1}}

% Note du transcripteur, sur l'état matériel du feuillet.
\newcommand{\note}[1]{\par\smallskip{\small\color{gray!60!black}\textit{[#1]}}\par\smallskip}

% Note marginale de Grothendieck, distincte du corps du texte.
\newcommand{\marginal}[1]{\par\smallskip\noindent\hspace{1em}%
  {\small\color{gray!60!black}#1}\par\smallskip}

% --- Titre ------------------------------------------------------------------

\AtBeginDocument{%
  \begin{center}
    {\large\bfseries Fonds Alexandre Grothendieck --- cote n\textsuperscript{o}~\grothFolder}\\[2pt]
    {\itshape\grothFoldertitle}\\[4pt]
    {\small feuillets \grothFirst--\grothLast\ (lot \grothBatch)}\\[2pt]
    {\footnotesize\color{gray!60!black}Transcription établie d'après le fac-similé numérisé
     par l'Université de Montpellier.\\
     Les lectures incertaines sont soulignées, les illisibles notées [\,\ldots\,],
     les ajouts entre crochets.}
  \end{center}
  \vspace{1em}}

\endinput


\folder{159}
\batch{2}
\pages{21}{33}
\dating{s.d. — le groupe « Dérivateurs » (157-1 à 160) est daté 1990-[1991]}
\watermark{Édition de démonstration}
\foldertitle{Connexions projectives et uniformisation : notes manuscrites (s.d.)}

\begin{document}

\page{21}
$\det E \simeq \det F \otimes L^{\otimes 2}$, donc $\underline{\omega}_X \simeq M \otimes L^{\otimes 2}$
\note{sous $\det E$ et sous $\det F$, reliés à eux par un double trait, il écrit
respectivement $\underline{\omega}_X$ et $M$}

donc (\uncertain{lorsque} $M$ est de degré $0$) $(2\deg L = \deg \underline{\omega}_X)$,
i.e.\ \struck{deg} $\underline{\deg L = g-1}$.

Considérons l'homomorphisme composé
\[
  \underline{\omega} \hookrightarrow E \simeq F \otimes L \twoheadrightarrow L
\]
c'est une section de $L \otimes \omega^{-1}$, qui est de degré $-(g-1) < 0$, donc
\ill{} est nulle, donc \struck{cela se factorise par} $E \to L$ se factorise par
$E/\underline{\omega} \simeq \underline{O}_X$, donc $L \simeq \underline{O}_X$,
absurde car $\deg L \neq 0$ !

Ainsi, le \struck{faisceau} \add{groupe} structural des \ill{}-fibrés projectifs
associés à une \ill{} ne peut se réduire à un sous-groupe de Borel.
\note{un symbole isolé, qui revient plus bas devant « fibré principal », n'a pas
été lu ; il qualifie les fibrés dont il est question}
Mieux, qu'il ne peut non plus se réduire à un sous-groupe
\uncertain{normalisateur} d'un tore maximal. En effet, si c'était le cas, alors
par \uncertain{passage au revêtement double}, on se \uncertain{restreindrait} à
un tore maximal, donc à un Borel, absurde (car le genre du revêtement
\uncertain{sera} aussi $\geqslant 2$ !). Le \ill{} argument montre :

\begin{quote}
\textbf{Th.} Soit $X$ courbe projective \add{lisse} de genre $g \geqslant 2$ sur
un corps $k$ (de car.\ nulle ?). Alors pour \textbf{tt} connexion projective sur
$X$, on ne peut réduire le \add{gpe} structural du \ill{}-fibré principal de
groupe \uncertain{$\mathrm{PGL}(1)_S$} $= H$ associé, à un sous-groupe
\struck{\ill{}} $H'$ distinct de $H$ \ill{}
\note{le nom du groupe est écrit d'un seul jet et les lettres n'en sont pas
assurées ; la lecture retenue est celle qu'impose la notation $\mathrm{PGL}(r)_S$
du début du dossier, où $r$ est la dimension relative du fibré de
Brauer-Séveri}
\end{quote}

\page{22}
Quand $k = \mathbb{C}$, cela signifie que l'hom.\ correspondant
\[
  \pi_1(X^{\mathrm{an}}, x) \longrightarrow H(\mathbb{C})
\]
a une image Zariski-dense dans $H_{\mathbb{C}}$.

\page{24}
$h\,\hat{c} = \hat{c}'$
\note{isolé et souligné en tête de page ; la suite reprend l'argument depuis
le début}

Une accolade réunit
\[
  (P,\sigma) \quad\longleftrightarrow\quad 0 \to \Omega \to E \to \underline{O}_S \to 0
\]
et $\widehat{S}$.

$P$ muni d'une \textbf{connexion $C_\rho$}. On cherche les
$P \xrightarrow{\;u\;} \widehat{S}$ (s'il y en a) pour lesquels $u(C_\rho)$ est
\struck{dans} de la forme $\hat{c}$.

La donnée d'un iso.\ à l'ordre $2$ donne \uncertain{déjà} l'extension
$0 \to \Omega \to E \to \underline{O}_S \to 0$ ; d'où $(P,\sigma)$ à
\uncertain{isom.} unique près.

Si on se donne une \ill{}
\note{la fin de la ligne, et une reprise interlinéaire qui paraît dire que les
$P$ partagent la connexion à l'ordre $2$, sont biffées d'un long trait}
$\widehat{S}$ conn.\ $\hat{c}$ : $C \in \mathrm{Conn}(\widehat{S})$. On cherche
alors les $g \in G$ (autom.\ de $\widehat{S}$ qui sont l'identité
\uncertain{mod} l'ordre $2$) tels que $g.C$ de la forme $\hat{c}$.
\marginal{Il existe tjs \ill{}}

\struck{\ill{}} on voit que $gC = k\,\hat{c}_0$, $(k^{-1}g)C = \hat{c}_0$.
\note{trois lignes reprises et biffées ; ce qui suit repart de $g_0$}

Quels sont les autres ? Ils sont de la forme $g = h\,g_0$, $(h \in G)$, et
\ill{} $gC$ $(= h(g_0 C) = h\,\hat{c})$ sera de la forme $\hat{c}'$, i.e.\
$h \in G(c)$. L'ens.\ de ces $h$ est en correspondance biunivoque avec
$G/G_0$ --- on peut aussi dire ainsi que pour \textbf{tt} extension de \ill{}
\note{la page s'arrête au bas du feuillet, la phrase inachevée}

\page{25}
\textbf{Th.} Soit $S$ un \uncertain{schéma} de car.\ $0$, $X/S$ un schéma,
$P = \varprojlim P_n$ une \struck{algèbre} $\underline{O}_X$-Algèbre
$\underline{O}_X$-augmentée adique complète, \struck{\ill{}} à filtration
$J$-adique, \add{(\ill{})} avec $J/J^2$ \struck{\ill{}} \ill{} inversible
\note{une reprise interlinéaire, biffée, commence par « tjs »}
(par une \ill{} l'augmentation),
\[
  J/J^2 = \omega \;\xrightarrow{\ \omega\ }\; \Omega^1_{X/S}
\]
\note{la ligne se poursuit par $\widehat{S} \otimes \ill{}\,/\,J\,\ill{}$,
en partie biffé et illisible}
\uncertain{$P$} muni d'une \struck{\ill{}} connexion $C$ rel.\ à $S$,
\uncertain{compatible} à la structure d'Algèbre \ill{}.

On dira \ill{} $P$ sur $X$, muni d'une section $\sigma$, (dans \ill{} la
\uncertain{catégorie} de \ill{} les $\mathcal{E}$ de
$0 \to \Omega \to E \to \underline{O}_S \to 0$), muni d'une connexion $c$ (en
respectant \uncertain{par ailleurs} la section \uncertain{unique} $\sigma$),
\struck{\ill{}}

\ill{} iso.\ d'algèbres augmentées
\[
  \widehat{S} \;\underset{\sim}{\overset{u}{\longleftrightarrow}}\; P
\]
\struck{tel qu} (où $\widehat{S} = \underline{O}_{\widehat{P}}$, $P$ \ill{}
$\sigma$ et $c$ ---) tel que $u(C) = \hat{c}$.

\begin{itemize}
  \item[a)] La donnée d'un iso.\ $u$ à l'ordre $2$ détermine $(P,\sigma)$,
  i.e.\ $\mathcal{E}$, à \uncertain{isom.}\ unique \uncertain{près}, par
  $\Omega = J/J^2$, $E = \Omega \otimes (J/J^3)$.
  \note{le signe entre $\Omega$ et la parenthèse est lu $\otimes$, sans
  certitude}
  \item[b)] $(P,\sigma)$ étant ainsi déterminé, pour \textbf{tt} isom.\ $u_3$
  à l'ordre $3$ près
  \[
    \widehat{S}/\widehat{J}^{\,4}_P \longrightarrow P/J^4
  \]
  prolongeant l'iso.\ à l'ordre $2$ \ill{}, il existe un \textbf{isomorphisme
  \ill{} un seul} $u$ qui le prolonge, et tel que $u(C)$ soit de la forme
  $\hat{c}$, et pour \ill{}
\end{itemize}
\note{le feuillet s'arrête ici ; b) reste inachevé}

\page{26}
$c$ (sur $P$) il existe un iso.\ et un seul $u$ tel
\struck{que $u(\hat{c})$ prolonge} $u_2$, tel que $u(\hat{c}) = C$.

D'où corr.\ biunivoque entre
\begin{itemize}
  \item[a)] les solutions des $P_3$ \uncertain{initial}, à isom.\ près
  (N.B.\ la cat.\ des solutions est \textbf{rigide}) ;
  \item[b)] $(P,\sigma)$ étant déjà connu, comme \ill{}, avec $u_2$ ---
  l'ens.\ des \add{pseudo-}$u_3$ qui prolongent $u_2$ --- qui forme donc un
  torseur \ill{} $\Gamma(X, \Omega^{\otimes 2})$ ;
  \item[c)] l'ens.\ des connexions $C$ sur $P$ telles que $\hat{c}$ soit,
  dans \ill{}, \uncertain{donnée} de $C$ \ill{} \uncertain{sur}
  $\mathrm{Conn}^{*}(P)$, \struck{\ill{}}
\end{itemize}

\note{le haut de la page s'arrête là ; ce qui suit, séparé par un blanc, est un
brouillon de calcul}

$\mathcal{D}_\omega$ : \ill{}

\begin{tikzcd}
  \mathrm{Sol} \arrow[r, "\sim"] \arrow[dr] & \mathrm{Iso}_3 \\
   & \mathrm{Conn}^{*}P
\end{tikzcd}

\note{les deux buts portent chacun un point d'exclamation ; la flèche oblique
porte un label non lu}

\[
  \Omega \longrightarrow \widehat{S} \otimes \Omega^1_{X/S},
  \qquad
  \prod_{i \geqslant 0} \Omega^{\otimes i} \otimes \Omega^1_{X/S},
  \qquad g \in G(c)
\]

$\Omega + \mathcal{O} + \Omega$ \quad $g\hat{c}_0$ \quad $\hat{c} = \hat{c}$

\struck{$g\,u(C) = \ill{}$} \quad $gC = \hat{c}$ \quad $u(C) = \hat{c}$ \quad
\struck{$g\hat{c} = \hat{c}$} \quad \struck{$g\hat{c} = C$} \quad
$g^{-1}\hat{c} = C$

\struck{$g\,u(C) = g\hat{c} = \hat{c}'$} \quad $u(C_\rho) = C = \hat{c}$ \quad
$g = \mathrm{id}$ \quad \struck{$gC = \hat{c}$} \quad $gC = k\,\hat{c}$ \quad
$h\hat{c} = \hat{c} \Rightarrow h = 1$ \quad $k^{-1}g\,C = \hat{c}$ \quad
$g \in G(c)$ \quad $g'C = \hat{c}'$

\note{ces égalités sont jetées côte à côte sur trois lignes, sans phrase ; la
mise en colonnes est de l'éditeur}

\page{27}
\[
  S = \mathrm{Sym}(\Omega) \;\xrightarrow{\ c\ }\; \widehat{S} \,\widehat{\otimes}\, \Omega^1_{X/S},
  \qquad
  \Omega \;\xrightarrow{\ c\ }\; \widehat{S} \,\widehat{\otimes}\, \Omega^1_{X/S}
\]
\[
  C(\omega_0) = \sum_{i \geqslant 0} \omega_0^{\,i} \otimes \varpi_i,
  \qquad \varpi_i \in \Omega^1_{X/S}
\]
\note{le second membre du produit tensoriel est écrit d'un seul trait bouclé,
lu $\varpi_i$ ici et dans toute la suite}

\[
  C(\lambda\omega_0) = \lambda\,C(\omega_0) + \omega_0\,d\lambda
\]
\[
  = 1 \otimes \lambda\varpi_0 + \omega_0 \otimes (\lambda\varpi_1 + d\lambda)
  + \omega_0^2 \otimes \lambda\varpi_2 + \cdots + \omega_0^{\,i} \otimes \lambda\varpi_i \cdots
\]

$\Omega$ : $J^2 \longrightarrow J \otimes \Omega^1_{X/S}$, \quad
$\Omega \longrightarrow \Omega^1_{X/S} = \Omega$.

\[
  d(\lambda\omega) = \lambda\,d\omega + \struck{\omega\,d\lambda}
\]
\note{le second terme est biffé et annoté « nul mod $J \otimes \Omega^1_{X/S}$ »}

$\Omega^{-1} \otimes \Omega^1_{X/S}$

\[
  J^i \longrightarrow J^{i-1} \otimes \Omega^1_{X/S}
\]
\[
  \lambda\omega^i \longrightarrow \struck{C(\ill{})}\ \omega^i \otimes d\lambda
  + \lambda\,\omega^{i-1} C(\omega)
\]
\[
  \struck{\Omega}\ \Omega^i \xrightarrow{\ \uncertain{h_i}\ } \Omega^{i-1} \otimes \Omega^1_{X/S}
\]
\[
  \lambda\omega^i \longrightarrow \lambda\,\omega^{i-1}\,\underline{C(\omega)},
  \qquad
  \lambda\omega_0^{\,i} \longrightarrow \lambda\,\omega_0^{\,i-1}\,C(\omega_0)
\]
$C(\omega) = \varpi_0\,\omega + \ill{}$, \quad $\Omega^1_{X/S} \otimes \Omega^{-1}$.

$\lambda\omega^i\,\varpi_0$ \note{entouré d'un trait} \quad $\mathrm{Conn}^{*}$
\quad $\hat{c}\,\lambda = \ill{}$
\note{le bas du feuillet ne porte plus que ces bribes, en crayon très pâle}

\page{28}
\[
  C(\omega_0^{\,u})^{u^{-1}} =
  \left(\varpi_0 + \frac{\omega_0}{1+\sum \alpha_i \omega_0^{\,i}}\,\varpi_1
  + \left(\frac{\omega_0}{1+\sum \alpha_i \omega_0^{\,i}}\right)^{\!2}\varpi_2\right)
  \left(1 + \sum (i+1)\alpha_i \left(\frac{\ill{}}{1 + \ill{}}\right)\right)
\]
\[
  + \sum \left(\frac{\omega_0}{\sum 1 + \alpha_j \omega_0^{\,j}}\right)^{\!\ill{}} \otimes\, d\alpha_i
\]
\note{il écrit le facteur en exposant : la page 30 pose en toutes lettres
$C(\omega_0^{\,u}) = C\bigl(\omega_0(1+\sum_{i\geqslant 2}\alpha_i\omega_0^{\,i})\bigr)$,
de sorte que $\omega_0^{\,u}$ désigne le produit $\omega_0\,u$ avec
$u = 1 + \sum \alpha_i \omega_0^{\,i}$, celui-là même qui figure aux
dénominateurs ; la notation est conservée telle quelle ci-dessous}

Supposons $\alpha_2 = 0$, et calculons mod $\Omega^4$ :
\[
  C(\omega_0^{\,u})^{u^{-1}} \equiv (\varpi_0 + \omega_0\varpi_1 + \omega_0^2\varpi_2)
  \left[1 + \struck{\ill{}}\ 4\alpha_3\omega_0^3\right]
\]
\[
  \sum_{i \geqslant 1} \Omega^i \otimes \Omega^1_{X/S}
  \qquad = \struck{\ill{}}\ C(\omega_0) + 4\alpha_3\,\omega_0^3\,\varpi_0
\]

Supposons $\alpha_2 = 0$, $\alpha_i = 0$ $(i > 3)$, \marginal{$\alpha_n = ?$}
calculons mod $\Omega^{n+1}$ :
\[
  C(\omega_0^{\,u})^{u^{-1}} = (\varpi_0 + \omega_0\varpi_1 + \omega_0^2\varpi_2)
  \bigl(1 + (n+1)\alpha_n\,\omega_0^{\,n}\bigr)
\]
\[
  = \struck{\varpi_0}\ C(\omega_0) + \varpi_0\ \struck{\ill{}}\ (n+1)\alpha_n\,\omega_0^{\,n}\,\varpi_0
\]
donc il faut $(n+1)\alpha_n = 0$.

Calculons \struck{mod} $\Omega^3$ en \uncertain{supposant} $\alpha_2$ \ill{} :
\[
  C(\omega_0^{\,u})^{u^{-1}} = (\varpi_0 + \omega_0\varpi_1 + \omega_0^2\varpi_2)
  \bigl(1 + 3\alpha_2\,\omega_0^2\bigr)
  = C(\omega_0) + 3\alpha_2\,\omega_0^2\,\varpi_0
\]

$\mathrm{Conn}(P)$,

\textbf{Prop.} Soit $\mathcal{C}$ l'ens.\ des connexions de $P$ (= de l'Alg.\
\add{de Lie} \ill{} $(J)$), \struck{\ill{}} $G_0$ le \struck{\ill{}} groupe des
automorphismes de $\widehat{S}$ (alg.\ augmentée) (qui sont l'identité mod
$J^4$), considérons l'application
\[
  (c,u) \longmapsto \struck{\ill{}}\ u(\hat{c})
\]
\[
  \mathrm{Conn}\,P \;(\times\, G_0) \longrightarrow \mathrm{Conn}_{\mathrm{Alg}}(\widehat{S})
  \quad \struck{mod}
\]
Cette \uncertain{application} est \textbf{bijective} (\ill{} de car.\ $0$).

\page{29}
\textbf{Cor.} Supposons $u \in G_0$, \struck{\ill{}} et $u\hat{c}$ de la forme
$\hat{c}'$, alors $u = \mathrm{id}$.

\ill{}, \struck{$u\hat{c}$} $u\cdot\hat{c} = 1\cdot\hat{c}'$ $\Longrightarrow$
$(u,\ill{}) = \mathrm{id}$.

Soit $G$ le gpe des automorphismes \uncertain{formels} de $\widehat{S}$ (qui sont
l'identité mod $J^3$) (donc $G/G_0 \simeq \Gamma(\Omega^{\otimes 2})$), et
considérons
\[
  G \times \mathrm{Conn}(P), \qquad (g,c) \longmapsto g\cdot\hat{c}.
\]
Quand a-t-on $g\hat{c} = g'\hat{c}'$, i.e.\ $(g'^{-1}g)\hat{c} = \hat{c}'$, i.e.\
quand $h\hat{c}$ est-il de la forme $\hat{c}'$ ? C'est \ill{}
\[
  h \in G(c) \iff h\cdot\hat{c} \ \text{de la forme}\ \hat{c}
\]
\note{le membre de droite est écrit en toutes lettres}
$h \in G(c)$, $g \in G$, $gh \in G(\ill{})$, $G(c)$ sous-groupe de $G$.

On trouve que \struck{$\Rightarrow g = 1$, donc \ill{}} $G \subset \ill{}$,
\[
  \struck{G(c) \times{}}\quad G_0 \times G(c) \longrightarrow G,
  \qquad (g,h) \longmapsto g\cdot h
\]
est bijectif : (si $gh = g'h'$, i.e.\ $(g'^{-1}g)\,h = h'$, $d'$ :
\struck{$(g'^{-1}g)$} $(g''h)\cdot\hat{c}$ de la forme $\hat{c}''$,
$g''(h\hat{c}) $ \note{$h\hat{c}$ est souligné et annoté $\hat{c}'$}, d'où
$g'' = 1$, d'où \struck{$g = g'$} et par suite $h = h'$).

\struck{\ill{}} \uncertain{Surjectif} : soit $g \in G$, $g\hat{c} = k\,\hat{c}_0$,
d'où $k^{-1}g = h \in G(c)$, d'où $g = kh$.

\page{30}
\[
  C(\omega^{\,u}) = \struck{\ill{}}\ C\omega + \sum_{i \geqslant 2} C(\omega)u_i
  + \omega\,C(u_i)
  = C(\omega)\Bigl(1 + \sum u_i\Bigr) + \omega \sum C(u_i)
\]

\struck{Preuve}

$C|_{\Omega}$ : $\Omega \longrightarrow \widehat{S} \otimes \Omega^1_{X/S}$, \quad $\Omega^{\otimes i}$, \quad
$C(\omega) = \omega\,\varpi_{-1} + \omega\,\varpi_0 + \ill{}$
\note{le premier terme de ce développement n'est pas lu avec certitude}

\[
  C(\lambda\,\omega^{\otimes i}) = \omega^{\otimes i} \otimes d\lambda
  + \lambda\,i\,\omega^{\otimes i-1}\,\ill{}
\]

$\omega = \lambda\omega_0$, \quad $u_i = \struck{\ill{}}\ \alpha_i\,\omega_0^{\,i}$.
\[
  C\bigl((\lambda\omega_0)^{\,u}\bigr) =
  \bigl(\lambda\,\add{C(\omega_0)} + \omega_0 \otimes d\lambda\bigr)
  \Bigl(1 + \sum_{i \geqslant 2}\alpha_i\,\omega_0^{\,i}\Bigr)
  + \lambda\omega_0 \sum_{i \geqslant 2}
  \bigl(\omega_0^{\,i} \otimes d\alpha_i + i\,\alpha_i\,\omega_0^{\,i-1}\,C(\omega_0)\bigr)
\]
\[
  C(\omega_0) = \struck{\ill{}}\ \varpi_0 + \omega_0\varpi_1 + \omega_0^2\varpi_2,
  \qquad \varpi_i \in \struck{\ill{}}\ \widehat{S} \otimes \Omega^1_{X/S}
\]
\[
  C(\struck{\lambda}\,\omega_0^{\,u}) = \struck{C(\omega_0}\ (\varpi_0 + \varpi_1 + \varpi_2)
  \Bigl(1 + \sum_{i \geqslant 2}\alpha_i\,\omega_0^{\,i} + \ill{}\Bigr)
\]
\[
  C(\omega_0^{\,u}) = C\Bigl(\omega_0\bigl(1 + \sum_{i \geqslant 2}\alpha_i\,\omega_0^{\,i}\bigr)\Bigr)
  = C(\omega_0)\Bigl(1 + \sum \alpha_i\,\omega_0^{\,i}\Bigr)
  + \omega_0 \sum_{i \geqslant 2}\bigl(\omega_0^{\,i} \otimes d\alpha_i
  + i\,\alpha_i\,\omega_0^{\,i-1}\,C(\omega_0)\bigr)
\]
\[
  = C(\omega_0)\Bigl(1 + \sum_{i \geqslant 2}(i+1)\alpha_i\,\omega_0^{\,i}\Bigr)
  + \omega_0^2 \sum_{i \geqslant 2} \ill{}
\]
\[
  C(\omega_0^{\,u})^{u^{-1}} = \struck{\ill{}}
  \Biggl(1 + \sum_{i \geqslant 2}(i+1)\alpha_i
  \biggl(\frac{\omega_0}{1 + \sum_j \alpha_j\,\omega_0^{\,j}}\biggr)^{\!i}\Biggr)
  + \ill{} \sum_{i \geqslant 2}
  \biggl(\frac{\omega_0}{1 + \sum_{j \geqslant 2} \alpha_j\,\omega_0^{\,j}}\biggr)^{\!i+1}
  \otimes\, d\alpha_i
\]
\struck{$C(\omega_0)^{u^{-1}} = \varpi_0 + \varpi_1 + \varpi_2 \ill{}$}
\note{la dernière ligne du feuillet est entièrement biffée}

\page{31}
$G \supset G_0$ : $G$ gpe, $G_0$ s-gpe.\ \struck{\ill{}}

$C \supset C_0$ : $C$ $G$-ens., $C_0$ \uncertain{sous-ens.}
\note{cette page est le lemme de groupes tout seul, détaché de la géométrie}

\[
  G_0 \times C_0 \;\xrightarrow{\ \sim\ }\; C, \qquad (h,c) \longmapsto h\cdot c
\]
\[
  G(c) = \{\, g \in G \ \mid\ g\cdot c \in C_0 \,\}
\]
Soit $c \in C_0$, et supposons \ill{} ; alors $G(c) \xrightarrow{\ \sim\ } G/G_0$,
i.e.\ \struck{\ill{}}
\[
  G_0 \times G(c) \;\xrightarrow{\ \sim\ }\; G, \qquad (h,g) \longmapsto h\,g
\]
\textbf{bij.}\ car : si $hg = h'g'$ ($h,h' \in G_0$, $g,g' \in G(c)$), alors
$(h'^{-1}h)\,g = g'$, d'où $h''(gc) = g'c = 1\cdot g'c$
\note{sous $gc$ et sous $g'c$ il note chaque fois $C_0$}
donc $h'' = 1$, i.e.\ $h = h'$ et $g = g'$ : d'où l'injectivité.

Surjectivité : soit $h \in G$, alors $h\,c \in C$ de la forme $h\cdot c'$
($h \in G_0$, $c' \in C_0$), donc $\underbrace{h^{-1}k}_{g}\,c = c'$,
$h^{-1}h = g \in G(c)$, donc $h = h\,g$, $h \in G_0$, $g \in G(c)$.

\textbf{Cor.} $C/G_0$ \ill{}

\uncertain{Soit} $g \in G$, $c \in C_0$ : $gc = c \Rightarrow g = 1$.
\struck{On a \ill{} $g \in G(c)$, $1 \in$} considérons une orbite $U$ de $G$ dans
$C$, soit $G\cdot c$, et considérons $U \cap C_0$. Alors
$U \cap C_0 \simeq \struck{\ill{}}\ G_0 \backslash G$. \struck{\ill{}}
$c = h\cdot c_0$ ($h \in G_0$, $c_0 \in C_0$, uniques) donc \ill{}

\page{32}
$X/S$ schéma relatif.

$P$ : $\underline{O}_X$-Algèbre augmentée, sép.\ complète pour \uncertain{$J$},
top.\ $J$-adique ($J$ idéal d'augmentation), $J/J^2 = \Omega$ inversible.
\[
  C : P \longrightarrow P \,\widehat{\otimes}\, \Omega^1_{X/S}
\]
connexion d'Algèbre \ill{} \struck{compatible} (compatible avec \ill{}
d'augmentation) :
\[
  \begin{cases}
    C(fg) = f\,C(g) + g\,C(f), & f, g \ \text{sections de}\ P \\
    C(\lambda \cdot 1) = 1 \otimes d\lambda, & \lambda \ \text{section de}\ \underline{O}_X
  \end{cases}
\]

\struck{\ill{}} \textbf{N.B.} Loc.\ (sur $X$, \ill{} affine) on écrit
\struck{pour} \uncertain{un iso.} $P \simeq \widehat{\mathrm{Sym}}(\Omega)$
(comp.\ \ill{} aux \ill{} augm.) et $C$ comme \ill{}, \ill{} relatif.

$C|\Omega$, qui \ill{} :
\[
  \Omega \longrightarrow \widehat{\mathrm{Sym}}(\Omega) \,\widehat{\otimes}\, \Omega^1_{X/S}
  = \prod_{i \geqslant 0} \mathrm{Sym}^i(\Omega) \otimes \ill{}
\]
\note{sous $\mathrm{Sym}^i(\Omega)$ il récrit $\Omega^{\otimes i}$}
$C_\Omega = (\varpi_i)$ : les $\varpi_i$ pour $i \neq 1$ sont
\textbf{linéaires}, et
\[
  \begin{cases}
    \varpi_i \in \struck{\ill{}}\ \Gamma\bigl(\Omega^{\otimes (i-1)} \otimes \Omega^1_{X/S}\bigr) & (i \neq 1) \\
    \varpi_1 \ \text{une connexion sur}\ \Omega
  \end{cases}
\]
(On voit sur ceci que \ill{}.)

$C(J^i) \subset J^{i-1} \,\widehat{\otimes}\, \Omega^1_{X/S}$, $C$ induit un hom.\
$(i \geqslant 1)$ donc
\[
  \mathrm{gr}^i(S) = J^i/J^{i+1} \simeq \Omega^i
  \;\xrightarrow{\ \widetilde{c}\ }\;
  \struck{\ill{}}\ \bigl(P \,\widehat{\otimes}\, \Omega^1_{X/S}\bigr)
\]
qui est de type $\underline{O}_S$/\textbf{linéaire}, et
$= \Omega^{i-1} \otimes \Omega^1_{X/S}$, $\mathrm{Sym}^i(\Omega)$-linéaire, donné
par
\[
  \widetilde{c}(\omega_0) = \omega_i \otimes \varpi_0
\]

\page{33}
\[
  \varpi_0 \in \Gamma\bigl(\Omega^{-1} \otimes \Omega^1_{X/S}\bigr)
  = \mathrm{Hom}\bigl(\Omega, \Omega^1_{X/S}\bigr)
\]
est indépendant de $i$, et est le $\varpi_0$ de \ill{} : l'homomorphisme \ill{}.

On dit que $C$ est \uncertain{connexion} \textbf{\uncertain{régulière}} si la
section $\varpi_0$ de $\Omega^{-1} \otimes \Omega^1_{X/S}$ est
\textbf{\uncertain{régulière}}
(\struck{\ill{}} telle que $\varpi_0 : \Omega \xrightarrow{\ \sim\ } \Omega^1_{X/S}$ !).

\textbf{Prop.} (\struck{Supposons} \add{(car.\ nulle)}) $P = \widehat{S} = \underline{O}_{\widehat{P}}$, \ill{} $(P,\sigma)$ \ill{}.

(Soit $\mathrm{Conn}^{*}(P)$ \note{il glose en dessous « régulières »} l'ens.\ des
conn.\ \ill{} de $P$, \struck{telles que \ill{}} $\hat{c}$, d'où
$c \longmapsto \hat{c}$,
\[
  \mathrm{Conn}(P) \longrightarrow \mathrm{Conn}(\widehat{S})
\]
et soit $\mathrm{Conn}^{*}(\widehat{S}) = \mathrm{Conn}(\widehat{S})$ \ill{} des
conn.\ \textbf{régulières} (de $\widehat{S}$), $\mathrm{Conn}^{*}(P)$ son
image inverse dans $\mathrm{Conn}(P)$.

\struck{Soit} $G$, $G_0$ les \struck{gpes} \ill{} auto.\ de $\widehat{S}$
augmentées qui sont l'identité mod $J^3$ ($J^4$).

\add{Alors} \struck{\ill{}} l'application
\[
  \mathrm{Conn}^{*}(P) \times G_0 \longrightarrow \mathrm{Conn}^{*}(\widehat{S}),
  \qquad (c,g) \longmapsto g\cdot\hat{c}
\]
\ill{} est bijective (\ill{} bijective).

Soit $c \in \mathrm{Conn}^{*}(P)$. \textbf{Cor.} (L'ens.\ $G(c) \subset G$ des
$g \in G$ tels que $g\hat{c} \in \mathrm{Im}\,\mathrm{Conn}^{*}(P)$ est tel que
\[
  G(c) \;\xrightarrow{\ \sim\ }\; G/G_0 \simeq \Gamma(X, \Omega^{\otimes 2}).
\]

\end{document}
